\section{Subjetividad: el significado de ingeniería de «ser uno mismo» en la era de la IA}

Uno de los conceptos más fuertes de la entrevista es «construir la subjetividad». Tristan dice que, en la era de la IA, lo único importante quizá sea «ser uno mismo», pero aquí «ser uno mismo» no es un mensaje de autoayuda, sino lograr que la IA sepa quién eres: tus valores, tu estética, tus activos históricos, tu forma de expresarte, tus objetivos actuales y tus límites. Cuanto más clara sea la subjetividad, mejor podrá el avatar actuar en tu nombre.

\begin{figure}[H]
\centering
\includegraphics[width=0.94\textwidth]{subjectivity-stack.png}
\caption{Construcción de la subjetividad: la pila de contexto que permite a la IA actuar en tu nombre.}
\end{figure}

\begin{knowledgebox}{Lectura de la figura: por qué la subjetividad es un problema de ingeniería}
En la figura, la subjetividad se compone de los activos históricos, los valores y la estética, los objetivos actuales, los límites de conducta y la expresión agentic. No es una personalidad abstracta, sino un conjunto de entradas que determina si el avatar de IA se parece a ti, si puede tomar decisiones por ti y si sabe cuándo detenerse.
\end{knowledgebox}

\subsection{Por qué «ser uno mismo» cobra importancia}

Cuando la IA puede escribir artículos, publicar comentarios, reconocer personas, generar contenido y recomendar conexiones por iniciativa propia, lo único insustituible del usuario puede pasar a ser «quién es». Si una persona ha acumulado durante mucho tiempo obras, opiniones, estética y relaciones, la IA puede amplificarla con mayor precisión; si una persona carece de subjetividad, la IA solo puede generar una expresión promediada.

\begin{importantbox}{La subjetividad es el activo personal del futuro}
En la era de internet, el tráfico y los seguidores eran activos; en la era de las redes sociales con IA, la subjetividad que un modelo puede comprender y representar se convertirá en un activo. Incluye obras, expresión, valores, relaciones y límites.
\end{importantbox}

\subsection{Los límites del avatar}

Cuanto más se parece el avatar al usuario, mayor es el riesgo: puede decir algo equivocado, expresarse en nombre del usuario sin autorización o filtrar información privada. Por eso la subjetividad debe incluir límites: qué contenido puede hacerse público, qué contenido requiere confirmación y qué contenido el avatar nunca puede usar.

\begin{warningbox}{Un avatar sin límites no es un producto, es un riesgo}
Si un avatar solo busca «parecerse a ti», pero carece de límites de permisos, de auditoría y de mecanismos de revocación, se convierte en un agente que el usuario no puede controlar. Los productos sociales con IA deben construir a la vez la semejanza y la sensación de control.
\end{warningbox}

\subsection{Resumen de la sección}

La subjetividad es la entrada central de las redes sociales con IA. Permite que el avatar actúe en nombre del usuario y, al mismo tiempo, exige que el producto ofrezca límites y mecanismos de control.
